% =====================================================================
%  ch03b.tex ------ 第 3 章 §3.2 Chern--Simons 泛函与 Ray--Singer 挠率
%  原书 p.111--p.116（PDF 131--136）
% =====================================================================

% [原文 p.111]
\section{Chern--Simons 泛函与 Ray--Singer 挠率}

\takeaway{在平坦联络 $\alpha$ 处，Chern--Simons 泛函的 Hessian 给出二次形式 $Q$。它在规范轨道横截方向上的正规化行列式满足 $\sqrt{\det d_\alpha^{*}d_\alpha}/\sqrt{|\det Q|}=T_\alpha^{1/2}$，其中 $T_\alpha$ 是 Ray--Singer 挠率。}

设 $M$ 为有向紧三维流形。这里设 $M$ 无边界，Lie 群 $G$ 取 $SU(2)$。设 $P$ 为 $M$ 上的 $G$ 主丛。以 $\mathcal{A}$ 表示 $P$ 的联络全体构成的空间，则 $\mathcal{A}$ 可与取值于 $G$ 的 Lie 代数 $\fg$ 的一次微分形式全体构成的空间 $\Omega^{1}(M,\fg)$ 等同。如 2.5 节所说明，Chern--Simons 泛函 $CS\colon\mathcal{A}\to\mathbb{R}$ 由
\[
  CS(A)=\frac{1}{8\pi^{2}}\int_{M}\Tr\left(A\wedge dA+\frac{2}{3}A\wedge A\wedge A\right),
  \qquad A\in\mathcal{A}
\]
定义。$P$ 的规范群 $\mathcal{G}$ 可与 $M$ 到 $G$ 的光滑映射全体 $\Map(M,G)$ 等同；在规范群 $\mathcal{G}$ 对 $\mathcal{A}$ 的作用
\[
  g^{*}A=g^{-1}Ag+g^{-1}dg,\qquad g\in\mathcal{G}
\]
之下，Chern--Simons 泛函变为 $CS(g^{*}A)=CS(A)+n$，$n\in\mathbb{Z}$。如 2.5 节所示，Chern--Simons 泛函的临界点是平坦联络。

% [Tutor Note: 这里的归一化系数 $1/8\pi^{2}$ 与 $2/3$ 是配套的：对闭的 $M$、$G=SU(2)$，$CS(g^{*}A)-CS(A)$ 恰取整数值 $n$，$n$ 是 $g\colon M\to SU(2)$ 的映射度。对一般的单 Lie 群，系数须按 Killing 型的长根归一化相应调整，才能保持 $n\in\mathbb{Z}$。]

其次描述规范群作用的无穷小版本。设规范群 $\mathcal{G}$ 的元素 $g$ 表为 $g=e^{tf}$ 的形式，则可写
\[
  g^{*}A=A+t(df+[A,f])+\mathcal{O}(t^{2}),\qquad A\in\mathcal{A}.
\]
因此，规范群作用的无穷小版本由
\begin{equation}
  d_{A}f=df+[A,f] \tag{3.19}
\end{equation}
所定的协变导数 $d_{A}$ 给出。令 $\alpha\in\mathcal{A}$ 为平坦联络，$\alpha$ 所附带的局部系统记作 $\fg_{\alpha}$。以 $\Omega^{j}_{\alpha}$ 表示取值于 $\fg_{\alpha}$ 的 $j$ 次微分形式全体，便得到以 $d_{\alpha}$ 为微分的 de Rham 复形
\[
  0\longrightarrow\Omega^{0}_{\alpha}\longrightarrow\Omega^{1}_{\alpha}
  \longrightarrow\Omega^{2}_{\alpha}\longrightarrow\Omega^{3}_{\alpha}\longrightarrow 0 .
\]
这里 $d_{\alpha}\circ d_{\alpha}=0$ 由 $\alpha$ 是平坦联络推出。

本节假定平坦联络 $\alpha$ 满足 $H^{*}(M,\fg_{\alpha})=0$。由 Poincaré 对偶性，本质的条件是 $H^{0}(M,\fg_{\alpha})=0$ 与 $H^{1}(M,\fg_{\alpha})=0$；而 $H^{0}(M,\fg_{\alpha})=0$ 对应于：由平坦联络 $\alpha$ 的和乐所得基本群 $\pi_{1}(M)$ 的线性表示不可约，

% [原文 p.112]
又，$H^{1}(M,\fg_{\alpha})=0$ 对应于 $\alpha$ 在平坦联络空间中孤立。

% [Tutor Note: 这里的 $\fg_{\alpha}$ 是“局部系统”，即以 $\pi_{1}(M)$ 的表示 $\rho_{\alpha}$ 为纤维、在 $M$ 上扭成的平坦向量丛的截面层；$H^{j}(M,\fg_{\alpha})$ 是局部系统数上同调（在 de Rham 意义下即用 $d_{\alpha}$ 定义的上同调）。$H^{0}=0$ 等价于 $\rho_{\alpha}$ 不可约，$H^{1}=0$ 等价于 $\alpha$ 在模去规范等价后的平坦联络模空间中是孤立点，这两条保证后面的算子行列式都有限。]

这里复习一下为本节讨论提供动机的经典渐近展开理论。首先从 Gauss 积分
\[
  \int_{-\infty}^{\infty}e^{-\mu x^{2}}\,dx=\sqrt{\frac{\pi}{\mu}},\qquad \mu>0
\]
出发。对它作解析延拓，得到
\begin{equation}
  \int_{-\infty}^{\infty}e^{i\lambda x^{2}}\,dx
  =\sqrt{\frac{\pi}{|\lambda|}}\;e^{\frac{\pi i}{4}\operatorname{sgn}\lambda},\qquad \lambda\in\mathbb{R}.
  \tag{3.20}
\end{equation}
设 $Q$ 是以 $x_{1},x_{2},\dots,x_{n}$ 为变量的非退化二次形式，$\operatorname{sgn} Q$ 为 $Q$ 的符号差。利用式 (3.20) 可得
\[
  \int_{\mathbb{R}^{n}}e^{iQ(x_{1},\dots,x_{n})}\,dx_{1}\cdots dx_{n}
  =\frac{\pi^{n/2}}{\sqrt{|\det Q|}}\;e^{\frac{\pi i}{4}\operatorname{sgn} Q}.
\]
其次，对实值函数 $f(x_{1},\dots,x_{n})$，考虑积分
\begin{equation}
  \int_{\mathbb{R}^{n}}e^{ikf(x_{1},\dots,x_{n})}\,dx_{1}\cdots dx_{n}
  \tag{3.21}
\end{equation}
当 $k\to\infty$ 时的渐近行为。设 $f$ 的临界点只有有限个，且都是 $\mathbb{R}^{n}$ 中的非退化点。先处理一元积分
\[
  g(k)=\int_{-\infty}^{\infty}e^{ikf(x)}\,dx .
\]
在 $f(x)$ 的临界点以外，各点对 $k\to\infty$ 时积分的贡献相互抵消，故只需计算来自临界点的贡献。为简单起见，设 $f(x)$ 只有一个临界点 $x_{0}$，且 $f''(x_{0})>0$。此时
\[
  g(k)\sim\int_{x_{0}-\epsilon}^{x_{0}+\epsilon}e^{ikf(x)}\,dx
  =\int_{u_{1}}^{u_{2}}e^{ik(f(x_{0})+u^{2})}\frac{2u}{f'(x)}\,du
\]
成立。这里作了变量替换 $f(x)-f(x_{0})=u^{2}$。其次，

% [原文 p.113]
利用
\[
  \lim_{x\to x_{0}}\frac{2u}{f'(x)}=\sqrt{\frac{2}{f''(x_{0})}},
\]
得到
\[
  g(k)\sim_{k\to\infty}\sqrt{\frac{2}{f''(x_{0})}}
  \int_{-\infty}^{\infty}e^{iku^{2}+ikf(x_{0})}\,dx .
\]
再由 Gauss 积分 (3.20) 可知，$g(k)$ 当 $k\to\infty$ 时渐近展开的主项为
\[
  \int_{-\infty}^{\infty}e^{ikf(x)}\,dx
  \sim\sqrt{\frac{2\pi}{k f''(x_{0})}}\;e^{ikf(x_{0})+\frac{\pi i}{4}} .
\]
一般地，在关于 $f(x_{1},\dots,x_{n})$ 的上述假设之下，积分 (3.21) 当 $k\to\infty$ 时渐近展开的主项可表为
\begin{equation}
  \int_{\mathbb{R}^{n}}e^{ikf(x_{1},\dots,x_{n})}\,dx_{1}\cdots dx_{n}
  \sim_{k\to\infty}
  \sum_{\alpha}\frac{\pi^{n/2}e^{ikf(\alpha)}}{k^{n/2}\sqrt{|\det H(f,\alpha)|}}
  \;e^{\frac{\pi i}{4}\operatorname{sgn} H(f,\alpha)} .
  \tag{3.22}
\end{equation}
这里求和遍历 $f$ 的临界点 $\alpha$ 全体，$H(f,\alpha)$ 表示 $f$ 在 $\alpha$ 处的 Hessian。

% [Tutor Note: (3.22) 的相位 $e^{\frac{\pi i}{4}\operatorname{sgn} H(f,\alpha)}$ 就是平稳相位法的 Morse 指标相位：$\operatorname{sgn} H$ 等于 Hessian 的负特征值个数减去正特征值个数，故每个临界点贡献 $e^{\pm i\pi/4}$。因子 $k^{-n/2}$ 说明主项相当于“单圈”贡献；一元情形 $f''(x_{0})>0$ 时相位为 $e^{i\pi/4}$，与 (3.20) 中 $\lambda>0$ 的情形一致。]

现在回到 Chern--Simons 泛函的情形。令 $\alpha\in\mathcal{A}$ 为平坦联络，$(\Omega^{*}_{\alpha},d_{\alpha})$ 为由 $\alpha$ 决定的 de Rham 复形。联络空间 $\mathcal{A}$ 是仿射空间，故 $\mathcal{A}$ 在 $\alpha$ 处的切空间 $T_{\alpha}\mathcal{A}$ 可与一次微分形式空间 $\Omega^{1}(M,\fg)$ 等同。切空间 $T_{\alpha}\mathcal{A}$ 上由
\[
  \langle A,B\rangle=-\frac{1}{8\pi^{2}}\int_{M}\Tr(A\wedge*B)
\]
定义了内积。关于这个内积，得到切空间的正交分解
\[
  T_{\alpha}\mathcal{A}=\Image d_{\alpha}\oplus\Ker d^{*}_{\alpha} .
\]
这里 $d_{\alpha}\colon\Omega^{0}_{\alpha}\to\Omega^{1}_{\alpha}$ 是式 (3.19) 所定义的协变微分，$d^{*}_{\alpha}\colon\Omega^{1}_{\alpha}\to\Omega^{0}_{\alpha}$ 是它的伴随算子。切空间的子空间 $\Image d_{\alpha}$ 可看作

% [原文 p.114]
规范群作用下的轨道 $\mathcal{G}\alpha$ 在 $\alpha$ 处的切空间。为计算 Chern--Simons 泛函在平坦联络 $\alpha$ 处的 Hessian，考察以 $t$ 为参数的联络族 $\alpha_{t}=\alpha+t\beta$ 下 $CS$ 的行为，可知
\[
  CS(\alpha_{t})=CS(\alpha)+\frac{t^{2}}{8\pi^{2}}\int_{M}\Tr(\beta\wedge d_{\alpha}\beta)
  +\mathcal{O}(t^{3})
\]
成立。因此，平坦联络 $\alpha$ 处的 Hessian 由二次形式
\[
  Q(\beta)=-\langle\beta,*d_{\alpha}\beta\rangle
\]
给出。Chern--Simons 泛函在无穷小规范变换下不变，故 $Q$ 沿规范群作用下的轨道 $\mathcal{G}\alpha$ 的方向退化，从而在商空间 $\Omega^{1}_{\alpha}/d_{\alpha}\Omega^{0}_{\alpha}$ 上定出非退化二次形式。又，$Q$ 可看作 $\Omega^{1}_{\alpha}/d_{\alpha}\Omega^{0}_{\alpha}$ 上的算子 $-*d_{\alpha}$。以下计算这个二次形式的行列式。它是无限阶行列式，其确切含义是由谱 zeta 函数定义的正规化行列式，这一点将在后文说明。首先，在
\[
  \Omega^{1}_{\alpha}\oplus\Omega^{3}_{\alpha}\cong\Omega^{1}_{\alpha}\oplus\Omega^{0}_{\alpha}
\]
上引入算子
\[
  P=\varepsilon(d_{\alpha}*+*d_{\alpha}) .
\]
这里 $\varepsilon$ 在 $\Omega^{0}_{\alpha}$ 上取 $\varepsilon=1$、在 $\Omega^{1}_{\alpha}$ 上取 $\varepsilon=-1$。把 $P$ 的作用关于直和分解
\[
  \Omega^{0}_{\alpha}\oplus\Image d_{\alpha}\oplus\Ker d^{*}_{\alpha}
\]
加以描述，可表为
\begin{equation}
  P=
  \begin{pmatrix}
    0 & -d^{*}_{\alpha} & 0\\
    -d_{\alpha} & 0 & 0\\
    0 & 0 & Q
  \end{pmatrix}
  \tag{3.23}
\end{equation}
的形式，并得到
\begin{equation}
  P^{2}=\Delta^{0}_{\alpha}\oplus\Delta^{1}_{\alpha} .
  \tag{3.24}
\end{equation}
这里 $\Delta^{j}_{\alpha}$ 是由 $d^{*}_{\alpha}d_{\alpha}+d_{\alpha}d^{*}_{\alpha}$ 给出的 $\Omega^{j}_{\alpha}$ 上的 Laplace 算子。

% [Tutor Note: 在三维流形上，$d^{*}_{\alpha}$ 与 $d_{\alpha}$ 通过 Hodge 星号相联系：在 $1$ 次形式上 $d^{*}_{\alpha}=-\ast d_{\alpha}\ast$（符号随次数与度规的号差而定）。正因如此，$Q=-*d_{\alpha}$ 才能写成 (3.23) 中的右上块，而 $P$ 是自伴算子，$P^{2}$ 才能按 Hodge 分解为两块 Laplace 算子的直和。$\varepsilon$ 的符号正是为了让 $P$ 自伴。]

% [原文 p.115]
一般地，Laplace 算子 $\Delta$ 的正规化行列式（regularized determinant）定义如下。把 $\Delta$ 的正特征值全体计及重数排列为
\[
  \lambda_{1}\le\lambda_{2}\le\cdots\le\lambda_{j}\le\cdots,
\]
并定义谱 zeta 函数为
\[
  \zeta_{\Delta}(s)=\sum_{j=1}^{\infty}\frac{1}{\lambda_{j}^{s}} .
\]
已知它在 $s$ 的实部充分大时为正则函数，并可解析延拓为全平面上的亚纯函数。特别地，$s=0$ 不是极点，$\zeta_{\Delta}(0)$、$\zeta'_{\Delta}(0)$ 的值有定义。把 $\Delta$ 的正规化行列式定义为
\[
  \det\Delta=\exp\bigl(-\zeta'_{\Delta}(0)\bigr).
\]
回到平坦联络 $\alpha$。Ray 与 Singer 证明了
\[
  T_{\alpha}(M)=\frac{(\det\Delta^{0}_{\alpha})^{3/2}}{(\det\Delta^{1}_{\alpha})^{1/2}}
\]
与流形的 Riemann 度量无关，是拓扑不变量。称 $T_{\alpha}$ 为 $M$ 关于平坦联络 $\alpha$ 的 Ray--Singer 挠率（Ray--Singer torsion）。对 $|\det Q|$ 赋予适当含义，便可证明下述命题。

\begin{proposition*}{3.4}
在平坦联络 $\alpha$ 处，Chern--Simons 泛函的 Hessian 经正规化后所得的行列式，与 Ray--Singer 挠率之间成立
\[
  \frac{\sqrt{\det d^{*}_{\alpha}d_{\alpha}}}{\sqrt{|\det Q|}}
  =T_{\alpha}^{1/2} .
\]
\end{proposition*}

\begin{proof}
由算子 $P^{2}$ 的表示 (3.24)，得到作为正规化行列式的等式
\[
  \det P^{2}=\det\Delta^{0}_{\alpha}\det\Delta^{1}_{\alpha} .
\]
因此绝对值 $|\det P|$ 有意义。又，由 $P$ 的 (3.23) 表示可知，$|\det Q|$ 是

% [原文 p.116]
按满足
\[
  |\det P|=\det(d^{*}_{\alpha}d_{\alpha})\,|\det Q|
\]
的方式定义才自然。因此
\[
  |\det Q|=\frac{(\det\Delta^{1}_{\alpha})^{1/2}}{(\det\Delta^{0}_{\alpha})^{1/2}}
\]
成立，与 $\Delta^{0}_{\alpha}=d^{*}_{\alpha}d_{\alpha}$ 合起来便证明了命题。
\end{proof}

% [Tutor Note: $\Delta^{0}_{\alpha}=d^{*}_{\alpha}d_{\alpha}$ 是因为 $\Omega^{0}_{\alpha}$ 上没有 $(-1)$ 次形式，$d_{\alpha}d^{*}_{\alpha}$ 这一项在 $0$ 次上恒为 $0$；指数 $3/2$ 与 $1/2$ 则来自 $\ast:\Omega^{0}\to\Omega^{3}$、$\ast:\Omega^{1}\to\Omega^{2}$ 的配对------在三维流形上 Hodge 星号把一个 $0$ 次形式变成 $3$ 次形式，故 $\Delta^{0}$ 与 $\Delta^{3}$ 的谱相同，这正是 $(\det\Delta^{0})^{3/2}$ 中 $3$ 的来源。]

关于 Chern--Simons 泛函的配分函数
\[
  Z_{k}(M)=\int\exp\bigl(2\pi\sqrt{-1}\,k\,CS(A)\bigr)\,\mathcal{D}A
\]
当 $k\to\infty$ 时的渐近行为，下面把从前面的考察中得到的启发叙述一下。Chern--Simons 泛函的临界点是平坦联络，为简单起见，假定它们只有有限个。由式 (3.22) 所述的有限维类比，可以期待渐近展开的主项来自平坦联络。Chern--Simons 泛函沿规范轨道退化，但上述命题中的 $\sqrt{\det d^{*}_{\alpha}d_{\alpha}}$ 可解释为规范轨道的体积元与 Lie 群的体积元之比。作为 Hessian 的正规化行列式出现的 Ray--Singer 挠率，与称为 Reidemeister 挠率的拓扑不变量一致；这就是 Cheeger 与 Müller 的定理。Witten 所预言的 $Z_{k}(M)$ 当 $k\to\infty$ 时的渐近行为具有
\[
  Z_{k}(M)\sim\frac{1}{2}e^{-3\pi i/4}
  \sum_{\alpha}\sqrt{T_{\alpha}(M)}\;
  e^{-2\pi i I_{\alpha}/4}\;
  e^{2\pi i(k+2)CS(A)}
\]
的形式。这里右端是关于平坦联络 $\alpha$ 的求和。$I_{\alpha}$ 是由 Atiyah--Patodi--Singer [16] 定义的、称为谱流（spectral flow）的量，是模 $8$ 确定的整数。渐近行为的描述含有关于流形 framing 与相位因子的微妙问题。Witten 的论文 [1] 中用 $\eta$ 不变量来叙述，而这里采用了 Atiyah [14] 的标准 2-framing 的叙述方式。

% [Tutor Note: 指数中的 $k+2$ 相当于把 level 平移 $k+h^{\vee}$，$h^{\vee}=2$ 正是 $SU(2)$ 的 dual Coxeter 数；这一平移与 2-framing 的选取有关。$I_{\alpha}$ 只模 $8$ 确定，而 $e^{-2\pi i(I_{\alpha}+8)/4}=e^{-2\pi iI_{\alpha}/4}$，故相位无歧义。]

% CHUNK ch03b DONE
